\documentclass[12pt]{article}
\usepackage[utf8]{inputenc}
\usepackage[T1]{fontenc}
\usepackage{hyperref}
\title{Suki Language Resources for K-12 Curriculum Development}
\author{Compiled by Claude Sonnet 5.5 (claude-sonnet-5-5) for Andr\'as Kornai}
\date{}
\begin{document}
\maketitle

\section{Introduction}
Suki (ISO 639-3 \texttt{sui}, also called Wiram; formerly Nausaku) is spoken by about 3,500 people in the Western Province of Papua New Guinea. It belongs to the Gogodala--Suki (Suki--Aramia River) group, where it is a family-level isolate. Literacy is low, and English is the language of formal schooling. Little material exists, and this report lists only what the searches turned up.

\section{Primary Educational Resources}
The searches found no Suki primers, textbooks or school programmes. No curriculum material in Suki could be identified.

\section{Linguistic and Reference Materials}
\begin{itemize}
\item Wikipedia, ``Suki language''. It summarises classification, speaker numbers and alternate names, and cites Hatfield (2008), a brief grammatical sketch, and Voorhoeve (1970) on the Suki--Gogodala subgroup. The cited works themselves were not located. \url{https://en.wikipedia.org/wiki/Suki_language}
\item Wikipedia, ``Gogodala--Suki languages''. \url{https://en.wikipedia.org/wiki/Gogodala%E2%80%93Suki_languages}
\item New Guinea World, ``Suki-Aramia River''. \url{https://newguineaworld.linguistik.uzh.ch/families/papuan-gulf/suki-aramia-river/start}
\item Ethnologue subgroup page. \url{https://ethnologue.com/subgroup/2275/}
\item WALS language page. \url{https://wals.info/languoid/lect/wals_code_suk}
\end{itemize}

\section{Religious and Community Resources}
\begin{itemize}
\item Joshua Project, Suki in Papua New Guinea. The search summaries mention a New Testament published in 1981, which comes from the Wikipedia article and was not confirmed on a scripture site. \url{https://joshuaproject.net/people_groups/15108/PP}
\end{itemize}

\section{Digital and Audio Resources}
No Suki audio, video, corpus or software resources were found.

\section{Curriculum Development Context}
Suki has few documented resources beyond a brief grammatical sketch and a 1981 New Testament. It is not known to be taught in schools, and English is the language of instruction. A native-language curriculum would have to start from new orthography, primer and materials development, working with the speaker community. Consulting the cited sketch and related Gogodala material would be a sensible first step.
\end{document}